\documentclass[11pt,a4paper]{article}

\usepackage[T1]{fontenc}
\usepackage[utf8]{inputenc}
\usepackage[margin=2.5cm]{geometry}
\usepackage{amsmath}
\usepackage{booktabs}
\usepackage{enumitem}
\usepackage{hyperref}

\hypersetup{
  colorlinks=true,
  linkcolor=blue,
  urlcolor=blue
}
\setlist{itemsep=0.35em, topsep=0.5em}

\title{Review and Improvement of the Temporal Validation Architecture\\[0.4em]
\large A Plain-Language Guide for LM-IDNet}
\author{}
\date{4 September 2026}

\begin{document}

\maketitle
\tableofcontents
\newpage

\section{Purpose of This Document}

LM-IDNet converts network traffic into consecutive ten-minute count windows.
Neighbouring windows are often related: a burst of traffic can continue for
more than ten minutes, a device can remain in the same operating state for
hours, and traffic can follow a daily routine. This creates an important
question:

\begin{quote}
Does the current dataset split give an honest measurement of how the detector
will behave on future traffic, or can temporal similarity make the reported
performance look better than it really is?
\end{quote}

This document reviews the current architecture, explains the relevant concepts
without assuming previous knowledge of time-series validation, reports checks
performed on the project data, and proposes a stronger validation design.

The main conclusion is:

\begin{quote}
The current split is a sound foundation because it is chronological and keeps
whole captures together. It should not be replaced by a random window split.
However, the project should add actual-timestamp overlap checks and grouped
rolling-origin validation before making a strong claim about performance on
future days. The existing locked final-test partition should remain untouched.
\end{quote}

\section{The Basic Unit: One Ten-Minute Window}

For each ten-minute interval $t$, LM-IDNet creates a count vector such as

\[
\mathbf{x}_t =
\left(
x_{t,\mathrm{TCP}},
x_{t,\mathrm{UDP}},
x_{t,\mathrm{SSDP}},
x_{t,\mathrm{ARP}}
\right).
\]

One row therefore describes how many packets of each category appeared in one
ten-minute period. The intervals are half-open:

\[
[t,t+10\text{ minutes}).
\]

This means a packet at exactly the starting time belongs to the window, while a
packet at exactly the ending time belongs to the next window. Consequently,
the windows produced from one continuous capture do not overlap.

Non-overlapping does not mean unrelated. Consider the following example:

\begin{verbatim}
10:00-10:10  cloud upload begins
10:10-10:20  cloud upload continues
10:20-10:30  cloud upload finishes
\end{verbatim}

The three rows cover different packets, but they were produced by one continuing
activity. Their counts can therefore be similar.

\section{Concepts Needed to Understand the Problem}

\subsection{Temporal Dependence and Autocorrelation}

Two observations are temporally dependent when knowing an earlier observation
provides information about a later one. In LM-IDNet, a high-traffic window may
make another high-traffic window immediately afterwards more likely.

\emph{Autocorrelation} is a numerical summary of this relationship. A lag-one
autocorrelation compares every window with the next window:

\[
\operatorname{corr}(x_t,x_{t+1}).
\]

For ten-minute windows, lag one means ten minutes later, lag six means one hour
later, and lag 144 means one day later. A value close to $1$ means that the
two measurements tend to rise and fall together. A value close to $0$ means
that a simple linear relationship is weak. A negative value means that high
values tend to be followed by low values, or vice versa.

Autocorrelation does not by itself prove that a model is wrong. It tells us that
the windows should not be treated as if they were independent repetitions of
the same experiment when calculating uncertainty or designing validation.

\subsection{What the Dirichlet--Multinomial Model Does and Does Not Model}

The Dirichlet--Multinomial model captures variation in the composition of TCP,
UDP, SSDP, and ARP traffic. It is useful when the category proportions vary
more than an ordinary Multinomial model expects.

The current model gives a probability to each window separately. In simplified
form, fitting several windows uses

\[
P(\mathbf{x}_1,\ldots,\mathbf{x}_R\mid\boldsymbol{\alpha})
=
\prod_{r=1}^{R}
P(\mathbf{x}_r\mid\boldsymbol{\alpha}).
\]

This treats the rows as exchangeable: their order does not change the fitted
model. The model does not contain a rule saying that the traffic in window
$t+1$ depends on the traffic in window $t$.

That is acceptable if the research goal is to learn the general distribution
of individual windows. It is not enough if the goal is to forecast the next
window from recent history or to claim that all window-level results are
statistically independent.

\subsection{Data Leakage}

Data leakage occurs when information that would not be available at deployment
time influences training or configuration choices. For example, a random split
could put a window from Monday morning in training and the immediately following
window in testing. Because those windows can be extremely similar, the test no
longer represents a genuinely unseen future period.

Leakage is more serious than ordinary temporal dependence. Dependence is a
property of the real traffic. Leakage is an experimental mistake that lets
future information influence the past.

\subsection{Pseudoreplication}

Pseudoreplication means counting related observations as if every one were a
fully independent experiment. Suppose a single continuing incident causes six
adjacent windows to be anomalous. Reporting this as six independent incidents
would exaggerate the amount of evidence.

Similarly, 144 windows from one day do not provide the same evidence as 144
windows drawn independently from 144 different days. The windows are still
useful, but confidence intervals and claims about repeatability should respect
their grouping by capture or day.

\subsection{Fitting, Calibration, Development Testing, and Final Testing}

LM-IDNet uses four chronological roles:

\begin{description}

\item[Fitting.] Estimate the Dirichlet parameter vector
\(\boldsymbol{\alpha}\), which describes normal traffic.

\item[Calibration.] Score separate known-normal windows and choose the lower
likelihood threshold. Calibration is part of model construction; it is not an
independent performance test.

\item[Development testing.] Compare possible design choices, find problems,
and improve the system. Results from this partition may influence the design.

\item[Final testing.] Measure the frozen system once after all design choices
have been made. Repeatedly inspecting the final set would gradually turn it
into another development set.

\end{description}

This separation is particularly important for anomaly detection. If the
threshold were chosen from the same attack-labelled data used to report recall,
the result would be optimistic.

\subsection{Time-Series Split and Rolling-Origin Validation}

A time-series split preserves the direction of time: earlier data train the
system and later data evaluate it. Rolling-origin validation repeats this idea
at several points in history. The end of the training period is called the
\emph{origin}; it rolls forward between folds.

An expanding version looks like this:

\begin{verbatim}
Fold 1: fit [Day 1 Day 2]       -> validate Day 3
Fold 2: fit [Day 1 Day 2 Day 3] -> validate Day 4
Fold 3: fit [Day 1 ... Day 4]   -> validate Day 5
\end{verbatim}

An expanding window retains all earlier data. A sliding window instead retains
only a fixed amount of recent history:

\begin{verbatim}
Fold 1: fit [Day 1 Day 2] -> validate Day 3
Fold 2: fit [Day 2 Day 3] -> validate Day 4
Fold 3: fit [Day 3 Day 4] -> validate Day 5
\end{verbatim}

Expanding windows answer, ``How well does the detector work as it accumulates
history?'' Sliding windows answer, ``How well does the detector work if old
behaviour is forgotten?'' Expanding validation is the appropriate first choice
for the current non-adaptive system. Sliding validation becomes more relevant
when the project studies behavioural drift and adaptation.

\subsection{Grouped Validation}

Grouped validation treats a complete capture or day as one indivisible unit.
All windows from that capture remain together. This prevents highly related
neighbours from being randomly distributed between fitting and evaluation.

For this project, the natural group is a source capture after timestamp overlap
has been resolved. A day is also useful for reporting because it corresponds to
a real future operating period.

\subsection{A Gap, Purge, or Embargo}

Some time-series evaluations leave an unused interval between training and
testing. This interval is variously called a gap, purge, or embargo. It is
necessary when observations overlap, when a feature uses future values, or when
the intended deployment question concerns performance after a period of time
rather than performance in the immediately following window.

An arbitrary gap is not automatically more scientific. If the detector will
score the very next ten-minute window in production, evaluating that next
window is realistic. The mandatory rule here is to eliminate actual duplicate
or overlapping evidence. A one- or two-hour gap can then be added as a
sensitivity experiment if training-only autocorrelation analysis shows that the
conclusion depends strongly on immediate neighbours.

\subsection{Block Bootstrap}

A bootstrap repeatedly resamples data to estimate how uncertain a metric is.
A normal row-by-row bootstrap would break the temporal groups and pretend that
all windows were independent. A block bootstrap resamples complete days or
captures instead. It therefore preserves much of the dependence inside each
block and gives more honest uncertainty intervals.

\section{Review of the Existing LM-IDNet Split}

\subsection{Current Structure}

The project configuration declares four lists of capture identifiers:

\begin{verbatim}
fit
calibration
development_test
final_test
\end{verbatim}

The validator requires every list to be chronological, prevents a capture
identifier from appearing twice, and requires each partition to end before the
next partition begins. Purpose-specific functions expose only the appropriate
capture list to training, calibration, development evaluation, or final
evaluation.

Relevant implementation files are:

\begin{itemize}

\item
\path{src/lm_idnet/config.py}, which validates the configured partition dates;

\item
\path{src/lm_idnet/partitioning.py}, which exposes immutable capture groups by
purpose; and

\item
\path{src/lm_idnet/processing/packet_transformer.py}, which produces UTC,
half-open, non-overlapping windows inside each capture.

\end{itemize}

\subsection{What Is Already Good}

The following architectural decisions should be retained:

\begin{enumerate}

\item The split is chronological rather than random.

\item A complete capture is assigned to one partition. Individual windows from
the same capture are not deliberately scattered among partitions.

\item Model fitting and threshold calibration use different captures.

\item Training and calibration cannot request the locked final-test identifiers
through their normal partition accessors.

\item The category order and window policy are centrally configured and
included in the dataset fingerprint.

\item Missing capture coverage is represented differently from an observed
silent window. Absence of packets is therefore not automatically confused with
absence of evidence.

\end{enumerate}

This is already much more meaningful than applying an ordinary random
train--test split to all ten-minute rows.

\subsection{Limitation: Dates Are Validated, Not the Actual Time Ranges}

The current configuration validator extracts a date from each capture name. It
does not compare the timestamps contained in the processed data. A filename can
therefore look correctly ordered even when its contents overlap another file.

The review found the following D-Link overlap:

\begin{itemize}

\item The nominal October 8 capture runs from 8 October at 08:10 UTC until
9 October at 16:00 UTC.

\item The nominal October 9 capture runs from 8 October at 16:00 UTC until
9 October at 16:00 UTC.

\item The files consequently share 144 ten-minute window identities.

\item Of those shared windows, 135 have identical count vectors and nine have
different count vectors.

\end{itemize}

Both files belong to the fitting partition. This is not leakage from testing
into training, but it gives the overlapping period approximately twice the
intended influence on the model. It is a form of duplicate evidence.

Three single-window timestamp overlaps were also found among Chromecast fitting
captures. Their counts differ, which is consistent with source files meeting
at a boundary or representing partial observations. They should still be
resolved explicitly rather than treated as independent windows.

No actual timestamp overlap was found between different configured partitions.
Thus the current accepted artifacts do not contain a detected fitting-to-test
timestamp leak, but the input uniqueness guarantee is incomplete.

\paragraph{Implemented resolution.}
The October 8 D-Link capture is now excluded from configured experiments. A
packet-level audit found repeated packet occurrences in that file, while the
October 9 capture is the cleaner source for their shared interval. The three
Chromecast boundary windows were found to contain complementary fragments with
no identical device packets across their paired source files. Their counts are
therefore summed through three explicit, fingerprinted configuration entries.
The pipeline now constructs this canonical view, rejects every unresolved
overlap, verifies actual partition time ranges, and writes a timeline-validation
report before model fitting.

\subsection{Limitation: One Chronological Split Gives One Version of History}

The current fixed split answers a useful but narrow question:

\begin{quote}
How did the model fitted on these particular earlier dates perform on these
particular later dates?
\end{quote}

It does not show whether the conclusion remains stable when the boundary moves
by one or two days. A threshold can work well on one future day and poorly on
another because of traffic drift, weekday effects, device state, or changes in
traffic volume.

The locked final set should not be used repeatedly to solve this problem.
Rolling-origin folds inside the development history provide repeated future
tests while preserving the final set for one unbiased evaluation.

\subsection{Limitation: Window Counts Can Be Mistaken for Independent Evidence}

The calibration minimum is currently expressed as a number of windows. A
one-day calibration set with 144 correlated windows may pass a requirement of
30 samples, even though it represents only one day of device behaviour.

Likewise, pooling every window into one confusion matrix is useful for counting
decisions, but it can hide large differences between days. A detector that
works on three days and fails on one should not be presented as uniformly
stable simply because the successful days contained more windows.

\section{Evidence of Temporal Dependence in the Current Data}

The following diagnostic correlations were calculated within individual
captures. Pairs were never created across capture boundaries. Values were
centred within each capture before the pairs were combined, which reduces the
effect of different daily averages.

\begin{table}[ht]
\centering
\caption{Examples of lag-one correlation in ten-minute observations.}
\begin{tabular}{lll}
\toprule
Dataset & Traffic measurement & Lag-one correlation \\
\midrule
D-Link Day Cam 5 & UDP, SSDP, and ARP counts & $0.46$ to $0.54$ \\
UNSW Chromecast & total packet count & $0.86$ \\
UNSW Samsung Camera & category proportions & $0.46$ to $0.60$ \\
\bottomrule
\end{tabular}
\end{table}

The current development anomaly scores also had lag-one correlations of about
$0.05$ for D-Link, $0.39$ for Chromecast, and $0.42$ for Samsung Camera.
This is encouraging for D-Link because the score is much less persistent than
several raw categories. Chromecast and Samsung still show meaningful clustering
of similar scores across neighbouring windows.

These values demonstrate dependence; they do not by themselves select a new
model or a mandatory gap size. That decision should be based on whether the
dependence causes unstable thresholds, clustered false alarms, or poor
future-day performance.

\section{A Small Rolling-Origin Diagnostic}

A read-only diagnostic was performed to see whether moving the fitting and
calibration boundary changes the result. Each fold used earlier captures for
fitting, the next capture for calibration, and the following presumed-benign
capture for validation. The configured one-percent lower-tail threshold was
used. The D-Link October 8 capture was omitted from this diagnostic because of
its overlap with October 9.

\begin{table}[ht]
\centering
\caption{Variation across preliminary chronological folds.}
\begin{tabular}{lcc}
\toprule
Dataset & Future benign windows flagged & Threshold range \\
\midrule
D-Link Day Cam 5 & $0.00\%$ to $2.78\%$ & $-27.22$ to $-20.63$ \\
UNSW Chromecast & $0.00\%$ to $13.19\%$ & $-27.87$ to $-16.79$ \\
UNSW Samsung Camera & one fold available & insufficient for a range \\
\bottomrule
\end{tabular}
\end{table}

A one-percent calibration quantile does not guarantee a one-percent false-alarm
rate on the next day. It guarantees only that the selected threshold lies near
the bottom one percent of the calibration scores. If the next day's traffic
distribution changes, its false-alarm rate can be much larger or smaller.

The Chromecast range shows why several chronological origins are valuable. The
current one-split result may be correct for its selected dates while still
failing to describe normal day-to-day variability.

This diagnostic is evidence for improving the protocol, not a final experiment.
The D-Link days are only presumed benign, the folds were not preregistered, and
Samsung has too little pre-attack history for repeated folds.

\section{Recommended Architecture}

The recommended design is a \emph{canonical unique timeline followed by
grouped, three-way, expanding rolling-origin validation and one locked final
test}.

\begin{verbatim}
Raw captures
    |
    v
Validate actual timestamps and construct one unique timeline
    |
    v
Group complete captures/days
    |
    +--> Fold 1: past fit -> later calibration -> next validation
    |
    +--> Fold 2: more past fit -> later calibration -> next validation
    |
    +--> Fold 3: more past fit -> later calibration -> next validation
    |
    v
Freeze every design choice
    |
    v
Evaluate the locked final period once
\end{verbatim}

\subsection{Stage 0: Construct a Canonical Unique Timeline}

Before statistical splitting, the project should establish one identity for
every device window:

\[
(\text{device},\text{window start},\text{window end}).
\]

The validation boundary should:

\begin{enumerate}

\item load the actual UTC start and end times from every processed capture;

\item reject any window identity that occurs in more than one source capture;

\item report the capture names, timestamps, counts, and whether the counts are
identical;

\item prevent model fitting until each overlap has a documented resolution;
and

\item verify again that no resolved window crosses two statistical partitions.

\end{enumerate}

The application should not silently discard one D-Link copy. The nine
disagreeing windows show that ``keep the first'' or ``keep the last'' would be
an unexplained scientific choice. The source captures should first be inspected
to decide whether one is authoritative, whether packets were duplicated, or
whether the files contain complementary partial traffic. The resolution and
its effect on the dataset fingerprint should be recorded.

\subsection{Stage 1: Keep Captures or Days Intact}

After overlap resolution, a complete capture or day should remain the atomic
validation group. Random row-level splitting should be forbidden for the main
anomaly-detection evaluation.

Random 10\%, 15\%, 20\%, and 25\% row samples remain appropriate in a separate
paper-reproduction experiment whose purpose is to measure estimator behaviour
as sample size changes. They should not be presented as chronological
validation of future detection.

\subsection{Stage 2: Use Three Ordered Blocks in Every Development Fold}

Ordinary time-series cross-validation commonly contains only training and
validation. LM-IDNet needs a middle calibration block because its anomaly
threshold is learned from data:

\begin{verbatim}
earlier normal captures
        -> fit alpha
later normal captures
        -> calibrate threshold
following captures
        -> measure future performance
\end{verbatim}

The same score definition, quantile, window duration, and other choices must be
used consistently within a fold. All fold definitions and random seeds should
be written before reading their outcomes.

\subsection{Stage 3: Report by Day Before Pooling}

Each validation capture should produce its own report containing:

\begin{itemize}

\item number of observed, silent, and missing windows;

\item fitted concentration and
\(\psi=1/\sum_k\alpha_k\);

\item calibration threshold and calibration-score distribution;

\item number and percentage of false alarms on known-normal windows;

\item precision, recall, and F1 score when attack labels exist;

\item number of contiguous anomaly episodes, not only the number of anomalous
windows; and

\item fitting and scoring runtime.

\end{itemize}

The summary should then report the median, minimum, and maximum across captures.
If uncertainty intervals are required, complete captures or days should be
resampled as blocks.

\subsection{Stage 4: Freeze the System and Use Final Data Once}

Rolling-origin results may be used to choose:

\begin{itemize}

\item raw or per-packet-normalised scoring;

\item calibration quantile;

\item window duration;

\item minimum amount of fitting and calibration history;

\item whether a temporal gap is needed; and

\item whether time-of-day conditioning is justified.

\end{itemize}

After these decisions, the configuration, code commit, dataset fingerprint,
and evaluation protocol must be frozen. The final-test captures are then scored
once. A later experiment may deliberately define a new protocol, but it must
not overwrite or replace the first final result without disclosure.

\section{Dataset-Specific Fold Suggestions}

These folds are examples to be reviewed after timestamp overlaps and benign
assumptions are resolved. They illustrate the design; they are not yet frozen
experimental specifications.

\subsection{D-Link Day Cam 5}

After resolving the October 8/9 overlap, two useful development folds are:

\begin{verbatim}
Fold 1
  Fit:         October 8-11
  Calibration: October 12-13
  Validation:  October 14-15

Fold 2
  Fit:         October 8-13
  Calibration: October 14-15
  Validation:  October 16 and 19

Locked final test
  October 21 and 22
\end{verbatim}

Because attack labels are not available for these D-Link captures, these folds
can support claims about model stability, score drift, and false-alarm behaviour
only if the validation days are credibly documented as benign. They cannot
support a measured attack recall claim.

\subsection{UNSW Chromecast}

The known-benign period permits two preliminary stability folds:

\begin{verbatim}
Fold 1
  Fit:         October 10-13
  Calibration: October 14-15
  Validation:  October 16-17

Fold 2
  Fit:         October 10-15
  Calibration: October 16-17
  Validation:  October 18-19

Attack-labelled development evaluation
  October 20-23

Locked final test
  October 25-27
\end{verbatim}

The benign folds measure threshold stability and false alarms. The labelled
development period measures attack detection. These are related but distinct
questions and should be reported separately.

\subsection{UNSW Samsung Camera}

Only four benign days precede the attack period. One preliminary chronological
check is possible:

\begin{verbatim}
Fit:         May 28-29
Calibration: May 30
Validation:  May 31
\end{verbatim}

The current attack-labelled development period is June 1--4 and the locked
final period is June 5--8. One benign fold is not enough to estimate variation
across origins. Samsung should therefore provide supporting evidence alongside
the other devices rather than being described as a repeated rolling-origin
study by itself.

\section{When a More Temporal Model Becomes Necessary}

Better splitting makes evaluation more honest; it does not change the model.
The current marginal Dirichlet--Multinomial detector may remain sufficient if
rolling-origin performance is stable and score autocorrelation does not cause
long clusters of false alarms.

If the improved evaluation exposes a problem, extensions should be introduced
one at a time and compared against the unchanged baseline.

\begin{description}

\item[Time-of-day profiles.] Fit a small number of profiles for broad periods,
such as night, morning, afternoon, and evening. This addresses a regular daily
schedule. Too many hourly profiles would leave too little data per model.

\item[Persistence rule.] Require two or more consecutive low scores before
opening an incident. This reduces isolated alerts, but it also delays detection
and must be evaluated explicitly.

\item[Sequential score monitor.] Apply a simple cumulative or exponentially
weighted monitor to the DM scores. This detects a sustained moderate change
that may not make any single window extreme.

\item[Sliding training window.] Use only recent known-safe history when fitting
a candidate model. This belongs to the later adaptation study because it raises
questions about drift, poisoning, promotion, and rollback.

\end{description}

None of these should be added merely because raw traffic is autocorrelated.
They are justified only when the grouped rolling-origin results demonstrate a
specific failure of the simpler detector.

\section{Recommended Implementation Order}

\begin{enumerate}

\item \textbf{Resolve duplicate time windows.} Investigate the D-Link October
8/9 overlap before producing another accepted model.

\item \textbf{Validate actual time ranges globally.} Add a fail-closed check for
duplicate window identities and cross-partition timestamp overlap.

\item \textbf{Define rolling folds as experiment data.} Store capture-level
fold definitions rather than computing random window splits.

\item \textbf{Run three-way development folds.} Fit, calibrate, and validate in
chronological order without touching final-test captures.

\item \textbf{Report results per capture.} Preserve daily values and then
summarise them; do not retain only one pooled metric.

\item \textbf{Estimate uncertainty by blocks.} Resample days or captures if
confidence intervals are needed.

\item \textbf{Freeze the experiment.} Record configuration, dataset
fingerprint, code commit, LM precision, and chosen settings.

\item \textbf{Run final evaluation once.} Use the explicit locked-final
boundary only after supervisor review of the frozen protocol.

\end{enumerate}

\section{Acceptance Checklist}

The improved architecture is ready when all of the following are true:

\begin{itemize}

\item Every processed device/window identity is globally unique or has a
documented, reproducible overlap resolution.

\item Actual timestamps, not only dates in filenames, are verified to move
forward across partitions.

\item No capture or day is divided among fitting, calibration, and evaluation.

\item Every rolling fold has fitting, calibration, and later validation blocks.

\item Fold definitions are fixed before their results are inspected.

\item Final-test captures are absent from all development folds.

\item Thresholds are recalibrated independently inside every fold.

\item Results are retained per capture or day as well as pooled.

\item Uncertainty calculations treat captures or days as blocks.

\item The experiment manifest records the configuration, dataset fingerprint,
code revision, precision, seeds, folds, and artifact hashes.

\item A separate, explicit command or gate is used for the one-time final
evaluation.

\end{itemize}

\section{What Can Be Claimed After This Improvement}

With the proposed architecture, the thesis can reasonably claim that:

\begin{itemize}

\item validation always predicts later traffic from earlier traffic;

\item neighbouring windows from one capture were not randomly divided between
training and testing;

\item duplicate timestamps were prevented from acting as repeated evidence;

\item threshold and detector stability were measured at several chronological
origins where the available data permitted it;

\item day-to-day variation remained visible instead of being hidden by one
pooled result; and

\item the final result came from an untouched period used after the design was
frozen.

\end{itemize}

It should still avoid claiming that every ten-minute window is independent, or
that a high number of windows automatically provides a high number of
independent experimental repetitions.

\section{Final Recommendation}

The current fixed chronological split should remain the outer research
protocol. It already protects the final period and is preferable to random
window sampling. The project should strengthen it with two additions:

\begin{enumerate}

\item a canonical, timestamp-unique dataset boundary; and

\item grouped three-way rolling-origin validation inside the development
history.

\end{enumerate}

This design is more meaningful because it measures the actual deployment
question---learning from the past and detecting changes in the future---at
several points in time, without sacrificing the honesty of the locked final
test. It also separates two different concerns correctly: rolling validation
makes the evaluation realistic, while a temporal model should be added only if
that realistic evaluation shows it is needed.

\end{document}
